\documentclass[11pt]{article}
\usepackage[a4paper,margin=25mm]{geometry}
\usepackage{amsmath,amssymb,amsthm}
\usepackage[T1]{fontenc}
\usepackage{hyperref}
\newtheorem{theorem}{Theorem}
\newtheorem{lemma}[theorem]{Lemma}
\newcommand{\AGL}{\operatorname{AGL}}
\newcommand{\orb}{\operatorname{orb}}
\title{A two-orbit maximal subgroup refutes source 00000001100}
\author{Local formalization report}
\date{8 October 2026}
\begin{document}
\maketitle

\begin{abstract}
The origin stabilizer in the affine group of a field line is a maximal proper
subgroup with exactly two orbits in the natural action. Over arbitrarily large
prime fields this contradicts the uniformly bounded error term asserted in
source 00000001100. The argument covers an eventual $q+O(1)$ interpretation as
well as the source's stronger ``always'' wording. The construction is standard
group theory; this report records a source-bound Lean formalization candidate.
\end{abstract}

\section{Statement and source scope}
The necessary first clause of the supplied English and Chinese source asserts
that the orbit count of a maximal subgroup of the finite-field affine group in
the natural action is $q$ plus a uniformly bounded constant. We take $q$ to be
the field cardinality, and maximality to mean maximality among proper subgroups
under inclusion. The one-dimensional affine group is an admissible instance of
that universal assertion.

The source additionally refers to a kernel of the constant and a subfield index,
without defining that clause. No replacement definition is supplied here. Since
the necessary first clause is false, its conjunction with any additional clause
is false. Nothing about the latter clause is required for this refutation.

For a field $F$, put $G=\AGL_1(F)$, the group of all affine bijections
$x\mapsto ax+b$ with $a\ne0$, acting on $F$ by evaluation. Let
\[
H=\{g\in G:g(0)=0\}.
\]
Orbit count means the cardinality of the quotient of $F$ by the actual
$H$-orbit equivalence relation.

\section{The subgroup and its orbits}
\begin{lemma}
The subgroup $H$ fixes $0$ and is transitive on $F\setminus\{0\}$.
\end{lemma}
\begin{proof}
Fixedness is its defining property. Given nonzero $x,y$, scaling by the unit
$y/x$ belongs to $H$ and sends $x$ to $y$. An affine bijection fixing zero cannot
send a nonzero point to zero, by injectivity.
\end{proof}

\begin{lemma}
$H$ is a maximal proper subgroup of $G$.
\end{lemma}
\begin{proof}
Translation by $1$ does not fix zero, so $H$ is proper. Suppose a subgroup
$J$ properly contains $H$. Take $f\in J\setminus H$, so $f(0)\ne0$. For any
$g\in G\setminus H$, transitivity on the nonzero points gives $h\in H$ with
$h(f(0))=g(0)$. Set $a=hf$, using composition multiplication. Then $a\in J$
and $a(0)=g(0)$, hence $a^{-1}g$ fixes zero and lies in $H\subseteq J$.
Therefore $g=a(a^{-1}g)\in J$. All elements of $H$ were already in $J$, so
$J=G$. This proves maximality.
\end{proof}

\begin{lemma}
The actual orbit quotient of $H$ on $F$ has cardinality $2$.
\end{lemma}
\begin{proof}
The first lemma shows that two points are orbit equivalent precisely when both
are zero or both are nonzero. Thus the quotient is in bijection with a two-element
set. Explicitly, send the zero class to \texttt{true} and the nonzero class to
\texttt{false}; the inverse selects the classes of $0$ and $1$. The field axiom
$0\ne1$ ensures that these are distinct. This applies also to the field of
order two.
\end{proof}

\section{An unbounded counterexample family}
\begin{theorem}
For every $C,q_0\in\mathbb N$, there is a prime $p\ge q_0$ and a maximal proper
subgroup $H_p$ of $\AGL_1(\mathbb F_p)$ such that
\[
\orb(H_p)=2,\qquad \orb(H_p)+C<p.
\]
\end{theorem}
\begin{proof}
By the infinitude of primes, choose a prime
$p\ge\max\{q_0,C+3\}$. Take $H_p$ to be the origin stabilizer. The previous
lemmas prove maximality and the count $2$, while $2+C<p$ follows from the choice
of $p$.
\end{proof}

Any eventual uniform $q+O(1)$ orbit-count assertion implies the necessary lower
bound $q\le\orb(H)+C$ for some fixed natural $C$ beyond some threshold $q_0$.
Indeed, if the source intended a real absolute error bound $B$, choose a natural
$C\ge\max\{B,0\}$, for example by taking a ceiling. For a signed error $c$,
$\orb(H)=q+c$ and $|c|\le B$ then imply $q\le\orb(H)+C$. The theorem refutes
this necessary inequality for every $C$ and every threshold. Consequently no
choice of positive, zero, or negative uniformly bounded error terms works.

\section{Formalization correspondence}
The Lean source uses mathlib's actual group \texttt{F $\simeq$\textsuperscript{a}[F] F}
of affine equivalences, and declares its evaluation action. The subgroup is the
actual \texttt{MulAction.stabilizer} of $0$; maximality is \texttt{IsCoatom} in
the subgroup inclusion lattice. Orbit count is
\begin{quote}\texttt{Nat.card (MulAction.orbitRel.Quotient H F)}\end{quote}
This is the actual action quotient, not a stipulated numerical function. A quotient equivalence with \texttt{Bool} proves its cardinality.

\texttt{arbitrarily\_large\_maximal\_counterexamples} quantifies both $C$ and
$q_0$ and constructs prime-field witnesses. \texttt{FiniteFieldOrbitCountClause}
quantifies all finite fields and all maximal subgroups, using
$q=\texttt{Nat.card F}$ and a signed integer deviation of uniformly bounded
absolute value. Restriction to prime fields yields a necessary consequence
already contradicted by the family. The final theorem is
\texttt{refutation\_00000001100 : $\neg$ FiniteFieldOrbitCountClause}.

The worker development compile uses Lean 4.33.0 and mathlib revision
\par\noindent\texttt{db584cd6d46c92f209a44c0f1c829460d327499d}.\par
The main declarations report
only \texttt{propext}, \texttt{Classical.choice}, and \texttt{Quot.sound}.
This is a worker proof candidate; independent semantic review and the manager's
deterministic package gate are separate requirements. Upstream submission
eligibility is not asserted here, and the known mathematical construction is
not claimed as a new discovery.
\end{document}
